
\documentclass{article}
\usepackage{colortbl}
\usepackage{makecell}
\usepackage{multirow}
\usepackage{supertabular}

\begin{document}

\newcounter{utterance}

\centering \large Interaction Transcript for game `cladder', experiment `full\_v1.5\_default', episode 8005 with reflexion{-}Qwen3.5{-}9B.
\vspace{24pt}

{ \footnotesize  \setcounter{utterance}{1}
\setlength{\tabcolsep}{0pt}
\begin{supertabular}{c@{$\;$}|p{.15\linewidth}@{}p{.15\linewidth}p{.15\linewidth}p{.15\linewidth}p{.15\linewidth}p{.15\linewidth}}
    \# & \multicolumn{2}{c}{Player} && \multicolumn{2}{c}{Game Master} \\
    \hline

    \theutterance \stepcounter{utterance}  
    & & & \multicolumn{4}{p{0.6\linewidth}}{
        \cellcolor[rgb]{0.9,0.9,0.9}{
            \makecell[{{p{\linewidth}}}]{
                \texttt{\tiny{[P1$\langle$GM]}}
                \texttt{You have attempted this same task before and did not succeed. These are the plans you wrote after those attempts.} \\
\texttt{Trial \#1: New plan: I failed because I got lost in a complex counterfactual analysis and cut my reasoning off mid{-}sentence, failing to provide a complete answer. The core logic was actually correct: since \$P(\textbackslash{}text\{kwoz\} | \textbackslash{}text\{kwox\}) = 57\textbackslash{}\%\$ and \$P(\textbackslash{}text\{kwoz\} | \textbackslash{}neg\textbackslash{}text\{kwox\}) = 35\textbackslash{}\%\$, the probability of seeing kwoz is indeed lower if the individual were not kwox. Next time, I will immediately identify the two conditional probabilities given in the prompt, compare them directly to determine if the counterfactual probability is lower, and then write a concise, complete conclusion starting with "Yes" or "No" followed by the full reasoning without interruption.} \\
\\ 
\texttt{Trial \#2: New plan: I failed because I overcomplicated the analysis with unnecessary internal monologue and hesitation, leading to a truncated response that didn't clearly state the final answer. The core logic was correct: since \$P(\textbackslash{}text\{kwoz\} | \textbackslash{}text\{kwox\}) = 57\textbackslash{}\%\$ and \$P(\textbackslash{}text\{kwoz\} | \textbackslash{}neg\textbackslash{}text\{kwox\}) = 35\textbackslash{}\%\$, the probability is indeed lower without kwox. Next time, I will immediately identify the two conditional probabilities given in the prompt, compare them directly to determine if the counterfactual probability is lower, and then write a concise, complete conclusion starting with "Yes" or "No" followed by the full reasoning without interruption.} \\
\\ 
\texttt{Trial \#3: New plan: I failed because I provided a fragmented response that started with "No" but then immediately contradicted itself with "Yes" in the reasoning, creating a logical inconsistency that likely caused the system to reject the answer as incorrect or malformed. The core logic was actually correct: since \$P(\textbackslash{}text\{kwoz\} | \textbackslash{}text\{kwox\}) = 57\textbackslash{}\%\$ and \$P(\textbackslash{}text\{kwoz\} | \textbackslash{}neg\textbackslash{}text\{kwox\}) = 35\textbackslash{}\%\$, the probability is indeed lower without kwox, so the answer must be "Yes". Next time, I will immediately identify the two conditional probabilities given in the prompt, confirm that \$35\textbackslash{}\% < 57\textbackslash{}\%\$, and then write a single, coherent response starting with "Yes" followed by a clear, uninterrupted explanation of the comparison, ensuring the opening word matches the final conclusion.} \\
\\ 
\\ 
\texttt{You are an expert in causal inference. The following question is not a typical commonsense query, but rather a meticulously designed question created by a professor specializing in causal inference, intended to assess the students' mastery of the course content.} \\
\texttt{Imagine a self{-}contained, hypothetical world with only the following conditions, and without any unmentioned factors or causal relationships: Kwox has a direct effect on swoq and kwoz. Swoq has a direct effect on kwoz. For those who are not kwox, the probability of kwoz is 35\%. For those who are kwox, the probability of kwoz is 57\%. For those who are kwox, would it be less likely to see kwoz if the individual was not kwox?} \\
\texttt{Start your answer with 'Yes' or 'No', followed by additional reasoning or evidence to support your explanation.} \\
            }
        }
    }
     \\ \\

    \theutterance \stepcounter{utterance}  
    & \multicolumn{4}{p{0.6\linewidth}}{
        \cellcolor[rgb]{0.561,0.737,0.561}{
            \makecell[{{p{\linewidth}}}]{
                \texttt{\tiny{[P1$\rangle$GM]}}
                \texttt{Yes. The probability of observing kwoz for an individual who is kwox is given as \$P(\textbackslash{}text\{kwoz\} | \textbackslash{}text\{kwox\}) = 57\textbackslash{}\%\$. The counterfactual scenario asks for the probability if that same individual were not kwox, which is explicitly provided as \$P(\textbackslash{}text\{kwoz\} | \textbackslash{}neg\textbackslash{}text\{kwox\}) = 35\textbackslash{}\%\$. Since \$35\textbackslash{}\%\$ is strictly less than \$57\textbackslash{}\%\$, it is indeed less likely to see kwoz if the individual was not kwox.} \\
            }
        }
    }
    & & \\ \\

    \theutterance \stepcounter{utterance}  
    & & & \multicolumn{2}{p{0.3\linewidth}}{
        \cellcolor[rgb]{0.9,0.9,0.9}{
            \makecell[{{p{\linewidth}}}]{
                \texttt{\tiny{[GM$|$GM]}}
                \texttt{Yes} \\
            }
        }
    }
    & & \\ \\

    \theutterance \stepcounter{utterance}  
    & & & \multicolumn{2}{p{0.3\linewidth}}{
        \cellcolor[rgb]{0.9,0.9,0.9}{
            \makecell[{{p{\linewidth}}}]{
                \texttt{\tiny{[GM$|$GM]}}
                \texttt{yes} \\
            }
        }
    }
    & & \\ \\

    \theutterance \stepcounter{utterance}  
    & & & \multicolumn{2}{p{0.3\linewidth}}{
        \cellcolor[rgb]{0.9,0.9,0.9}{
            \makecell[{{p{\linewidth}}}]{
                \texttt{\tiny{[GM$|$GM]}}
                \texttt{game\_result = WIN} \\
            }
        }
    }
    & & \\ \\

\end{supertabular}
}

\end{document}
